\section{Joint versus Single-Task Objectives}\label{sec:exp:ablation}

This section compares the joint loss of \cref{eq:method:loss_total} with training on only one of the two contrastive terms.
That is the comparison for RQ2.
Every run uses the Vanilla encoder of \cref{sec:method:arch}.
Single-task runs keep the dual-branch layout so that width and depth match.
The inactive branch is severed from the backward pass, so it does not affect the trained parameters and is not scored.

The three objectives are:
\begin{description}
  \item[Joint] $\Loss_{\mathrm{joint}} = \Loss_{\mathrm{em}} + \Loss_{\mathrm{md}}$, identical to Vanilla in \cref{sec:exp:results:vanilla}. Both branches are scored.
  \item[Emitter-only] $\Loss_{\mathrm{em}}$. The emitter branch is scored.
  \item[Mode-only] $\Loss_{\mathrm{md}}$. The mode branch is scored.
\end{description}

Joint numbers in \cref{tab:exp:ablation_em,tab:exp:ablation_md} are copied from the Vanilla rows of \cref{tab:exp:vanilla}, not from a second training run.
The single-task models use the same splits and the remaining hyperparameters of \cref{sec:exp:hparams}, except $\varepsilon$, which is selected independently on the validation grid (\cref{tab:exp:hparams_attn}).
The best value in each column is in bold and ties are bold on every tied row.

\begin{table}[tbp]
  \centering
  \caption{Emitter clustering for joint versus single-task objectives on the Vanilla encoder.}
  \label{tab:exp:ablation_em}
  \small
  \makebox[\linewidth][c]{%
  \begin{tabular}{@{}lccccc@{}}
    \toprule
    Objective & $\ARI$ & $\AMI$ & Hom. & Comp. & Count error \\
    \midrule
    Joint        & $0.948\pm0.176$ & $0.957\pm0.138$ & $0.989\pm0.052$ & $0.951\pm0.166$ & $0.25\pm0.49$ \\
    Emitter-only & $\mathbf{0.999\pm0.012}$ & $\mathbf{0.997\pm0.013}$ & $\mathbf{0.995\pm0.021}$ & $\mathbf{0.999\pm0.009}$ & $\mathbf{0.12\pm0.36}$ \\
    \bottomrule
  \end{tabular}%
  }
\end{table}

\begin{table}[tbp]
  \centering
  \caption{Mode clustering for joint versus single-task objectives on the Vanilla encoder. Entries follow the same aggregation as \cref{tab:exp:ablation_em}.}
  \label{tab:exp:ablation_md}
  \small
  \makebox[\linewidth][c]{%
  \begin{tabular}{@{}lccccc@{}}
    \toprule
    Objective & $\ARI$ & $\AMI$ & Hom. & Comp. & Count error \\
    \midrule
    Joint     & $\mathbf{0.985\pm0.107}$ & $\mathbf{0.986\pm0.102}$ & $\mathbf{0.994\pm0.027}$ & $\mathbf{0.988\pm0.103}$ & $\mathbf{0.23\pm0.50}$ \\
    Mode-only & $0.983\pm0.110$ & $0.984\pm0.103$ & $0.992\pm0.036$ & $\mathbf{0.988\pm0.102}$ & $0.24\pm0.56$ \\
    \bottomrule
  \end{tabular}%
  }
\end{table}

On emitters, Emitter-only outperforms Joint.
Mean $\ARI$ rises from $0.948$ to $0.999$, the spread collapses, and the count error falls to $0.12\pm0.36$.
Adding the mode term therefore costs deinterleaving performance on this backbone.

Mode-only is the other single-task ceiling.
Its mode $\ARI$ is $0.983$, slightly below Joint at $0.985$.
\Gls{umap} projections of two windows are in \cref{app:embeddings}.
An $\ARI$ summary of Feature \gls{dbscan}, Vanilla, Vanilla single-task, RoPE-TOA, and Bias is in \cref{app:ari_all}.

\Cref{fig:exp:kcorr_ablation_em,fig:exp:kcorr_ablation_md} show count errors as heatmaps of estimated versus true cluster counts.
Each cell is a window count.
The diagonal is exact recovery.
Mass above the diagonal is extra clusters (splits).
Mass below it is missing clusters (merges).
The Joint panels are the Vanilla evaluation of \cref{sec:exp:results:vanilla}, not a second run.
Only the active branch of each single-task run is shown.

On emitters, Emitter-only sits on the diagonal more tightly than Joint.
On modes, Joint is slightly higher on mean $\ARI$ than Mode-only.
Mode-only still sits with Joint on the heatmap, apart from a few windows that contain one true mode and are split into extra clusters.

\begin{figure}[tbp]
  \centering
  \expfig{kcorr_ablation_em.pdf}
  \caption{Estimated versus true number of emitters per test window for Joint and Emitter-only Vanilla. Cell values are window counts. The line is $\hat{K}=K$. The Joint panel is the Vanilla evaluation of \cref{tab:exp:vanilla}.}
  \label{fig:exp:kcorr_ablation_em}
\end{figure}

\begin{figure}[tbp]
  \centering
  \expfig{kcorr_ablation_md.pdf}
  \caption{Estimated versus true number of modes per test window for Joint and Mode-only Vanilla. Layout as in \cref{fig:exp:kcorr_ablation_em}.}
  \label{fig:exp:kcorr_ablation_md}
\end{figure}

RQ2, as stated, asks whether the joint objective outperforms a comparable transformer trained for deinterleaving alone.
On Vanilla emitters it does not.
On modes, Joint is slightly higher than Mode-only.
The next section asks whether physical structure in attention recovers the emitter accuracy of Emitter-only without dropping the mode term.
Emitter-only copies of RoPE-TOA and Bias are in \cref{sec:exp:attn_task}.
